\section{音楽嗜好を通じた自己開示}\label{音楽}
前節までで述べたように，自己開示は人間関係の親密化に重要な役割を果たすものの，現代の若者は深い自己開示を避ける傾向にある．また，深い自己開示に対しては，より慎重で共感的な反応が必要となる．この課題に対して，本研究では音楽嗜好を通じた自己開示に着目する．音楽嗜好は以下に挙げる4つの特徴を持つことから，自己開示を深めていく場合に適していると考えられる．


\subsubsection{パーソナリティとの関係}
Rentfrowら\cite{rentfrow_2003}\cite{rentfrow_2006}\cite{rentfrow_2007}の研究により，音楽嗜好とパーソナリティの間に有意な相関があることを実証的に示されている．具体的には，個人の音楽嗜好が，その人の性格特性，価値観，および認知的特徴を反映する傾向があることが明らかにされている．
このような音楽嗜好とパーソナリティの密接な関係性は，音楽を通じた自己開示が単なる趣味の共有以上の意味を持つことを示唆している．すなわち，音楽嗜好について語ることは，個人の内面的な特性を自然な形で表現する手段となる．


\subsubsection{普遍性}
音楽の普遍性は，対人コミュニケーションにおいて重要な意味を持つ．音楽は多くの人々が日常的に接している文化的要素であり，たとえ特定のアーティストや楽曲に対する知識の差があったとしても，共通の話題として取り上げやすいためである．つまり，音楽嗜好は会話のきっかけとして機能しやすく，自己開示を自然に促進する要素となる．


\subsubsection{楽曲に紐づく経験や感情}
音楽は，個人の経験や感情と密接に結びつく特性を持っている．この特性は，自己開示のプロセスにおいて重要な役割を果たす．具体的には，音楽に関する自己開示は，ジャンルや曲名といった表層的な情報の共有から始まり，特定の楽曲に対する個人的な思い入れや経験の共有へと，段階的に深化させることが可能である．例えば，特定の楽曲は，恋愛経験や人生の転機といった重要な個人的経験と結びついていることが多い．また，精神的な困難に直面した際の心理的支えとして機能することもある．


\subsubsection{音楽活動による絆の形成}
音楽活動による社会的な結びつきの形成について，Tarrら\cite{tarr_2014}は，社会的動機づけや快感，報酬に関わる神経系である内因性オピオイド系（EOS）と自己-他者の融合という2つのメカニズムに着目している．音楽の受動的な聴取でもEOSが活性化され，社会的な結びつきの形成が促進されることを示唆している．具体的には．音楽聴取がストレス軽減や気分の改善などの効果をもたらし，これらが人々の結びつきを促進する基盤として機能することが指摘されている．

\vskip\baselineskip

% 音楽嗜好を通じた自己開示がいいんすよ の話
これらの特徴から，現代の若者は音楽嗜好の共有をきっかけとして，その音楽にまつわる経験や個人的な価値観を話すことで，自然な流れで自己開示を行うことができる．このような段階的な相互開示のプロセスは，より深い相互理解と親密な関係性の構築へとつながると期待される．